\section{The walkthrough}
\label{part:walkthrough}

Ten steps take \eqref{eq:ni} from an informal wish to an empty ledger. Read the
margin panels as the primary text if you like; the prose exists to justify each
trade.

\clgreset

% ---------------------------------------------------------------- step 1
\begin{step}{What we are actually trying to prove}{%
  \clgnew{NI}{$\forall \coal.\ \Sec_\coal(P)$}%
  \clgnew{BIND}{$\Bind(\artifact,\pipeline)$}%
  \clgnote{BIND}{stays open almost to the end}%
}
\label{step:statement}

Open with the goal and with the one obligation that is easy to forget: a verdict
is a statement about \emph{a particular artifact}. ``This program is secure'' is
not a well-formed claim until you say which bytes you mean.

That second obligation is not pedantry. A confidentiality result computed on
LLVM IR captured before a late pass runs, whose mutated output then proceeds to
instruction selection, is a statement about a module that never executed. The
honest result in that situation is \Unknown{}, and nothing weaker than a binding
to the exact artifact --- bitcode hash, ordered pipeline digest, the capture
point, and the first IR-to-machine-IR boundary --- rules it out.

$\Sec_\coal$ is indexed by a \emph{coalition} $\coal$, a set of principals who
may pool what they observe. \cref{step:coalitions} explains why the index is
there. For now note that a single row in the panel stands for the whole family:
the ledger convention from \cref{sec:ledgerguide} is that an obligation
quantified over an index is one row with a subscript, not one row per instance.
\end{step}

% ---------------------------------------------------------------- step 2
\begin{step}{Declare the policy, and make the declaration total}{%
  \clgtouch{NI}%
  \clgnew{MAN}{$\Total(\manifest)$}%
}
\label{step:manifest}

Nothing about \eqref{eq:ni} is decidable until four questions have declared
answers: which inputs are secret, who observes what, which releases are
authorized and to whom, and which host executes each function.

These are supplied by a \emph{policy manifest}, and the obligation this step
opens is that the manifest is \emph{total} over everything reachable. Partiality
is the failure mode, not incorrectness: a reachable function with no declared
host has no observation projection, so no verdict follows for it --- and the
tempting default of treating an unplaced function as local and trusted is an
invented premise.

Two rules about the manifest are worth stating early because they are the ones
implementations break.

\textbf{Policy may not be inferred from names.} Not from an LLVM value name, not
from an argument position, not from \code{noalias}. This sounds like a style
guideline and is in fact load-bearing: it is what forces the manifest to be a
parsed input rather than a convention. It also has a blunt mechanical
justification --- \code{mlir-opt} discards SSA value names at parse, so a policy
encoded in a name does not survive to the analysis at all.

\textbf{Authority is not one relation.} A single principal-to-clearance map
cannot express the policies people actually want. The manifest instead carries
several independent relations:

\begin{center}
\begin{tabular}{@{}ll@{}}
\toprule
\emph{visibility} & who may \textbf{observe} an item; drives the projection \\
\emph{audience} & who a released value is \textbf{for}; drives the release premise \\
\emph{placement} & which host \textbf{executes} a function \\
\bottomrule
\end{tabular}
\end{center}

Visibility is itself carried by four separate bases --- over hosts, components,
outputs, and errors --- which is the concrete reason a single clearance level per
principal is insufficient: an item's observability is a function of where it
lives, not only of who is asking.

\begin{remark}[A relation the model does not yet have]
\label{rem:authorizers}
A natural fourth relation is \emph{who may authorize a release}, distinct from
who may read the released value. That separation is what would make ``may
declassify but may never read'' expressible --- an auditor trusted to decide
\emph{that} a summary may be published, without being trusted to see the raw
data behind it.

No such relation exists in the specification. Releases are identified by policy
and audience; there is no principal-level authorizer anywhere in the manifest
schema. So the auditor policy above is currently \textbf{not expressible}, and
the design examples that use an \code{authorizers} field are a proposed
extension rather than a description of the system. This is a genuine gap, not a
simplification for exposition.
\end{remark}
\end{step}

% ---------------------------------------------------------------- step 3
\begin{step}{Derive the coalitions; do not accept the authored list}{%
  \clgupd{NI}{$\forall \coal \in \coalitions.\ \Sec_\coal$}%
  \clgnew{CLO}{$\coalitions = \dnclose\coalitions_{\max}$}%
  \clgdone{CLO}%
  \clgnote{CLO}{immediate, and therefore easy to skip}%
}
\label{step:coalitions}

The coalition family is the \emph{downward closure} of the authored maximal
list, and it includes the empty coalition, which represents the world-visible
observer. Every member must be checked, including members nobody wrote down.

This step creates its obligation and discharges it immediately, because the
closure is a finite set operation. The reason to give it a step at all is that
``obvious and immediate'' is not the same as ``done,'' and the failure is silent:
an implementation that iterates the authored list rather than the derived closure
produces confident green results on artifacts that leak.

Two facts make the closure irreducible --- neither is an implementation detail.

\textbf{Verdicts are not monotone in the coalition.} The natural intuition is
that a larger coalition observes more, so checking the largest is the hardest
case and suffices. Release audiences break this. Authorization is
\emph{per-audience}, so whether a given release retires a pair's obligation
(\cref{sec:ledger}) depends on which coalition is asking. For a coalition inside
the declared audience the differing values are permitted and the obligation
retires; for a coalition outside it, nothing retires, the obligation is still
active, and \emph{the identical operation is a leak}. So \Verified{} at one
coalition and \Unsafe{} at another, with no containment between them.

\textbf{Joint visibility is not reconstructible from singletons.} An item may be
declared visible to $\{\textsf{alice},\textsf{bob}\}$ jointly and to neither
alone --- two shares of a secret, recoverable only together. A visibility
relation stored as principal-to-items cannot express this, and a report built by
evaluating singletons and taking unions cannot derive it.

Together these forbid deduplicating rows by coalition monotonicity, and forbid
omitting a derived coalition from the report. The cost is real and structural:
the closure over $n$ principals has $2^n$ members, each needing a row.
\end{step}

% ---------------------------------------------------------------- step 4
\begin{step}{Get inside a fragment you can actually reason about}{%
  \clgnew{NF}{$\NFConf(\artifact,\pipeline)$}%
}
\label{step:normalform}

Real LLVM IR is too large a language to have an exact semantics for. The move is
to fix a closed-world fragment --- a whitelist of types, opcodes, constants, and
intrinsics --- and to define the analysis only there. Everything outside is out
of profile.

The critical rule is what happens on the boundary. If an artifact does not
conform, the result is \Unknown{} with a reason naming the offending construct.
It is emphatically \textbf{not} a counterexample, even when the nonconforming IR
contains what looks unmistakably like a secret-dependent branch. A best-effort
pre-conformance scan may emit a triage finding, but that record is
non-authoritative and may never be presented as a replayable witness.

The reason for that asymmetry is that a counterexample is a claim about the
exact semantics, and outside the fragment there is no exact semantics to make a
claim about. Reporting \Unsafe{} there would be unsound in the same way that
reporting \Verified{} would be --- just less obviously so, and therefore more
dangerous, because a false alarm looks like diligence.

\begin{remark}[A scoping fact worth internalizing]
The fragment is scalar: integers and pointers, plus \code{float} and
\code{double}. Vectors are out of profile, as are the exotic float formats
(\code{half}, \code{bfloat}, and the 80- and 128-bit forms). So an analysis
defined on it says nothing about vector or tensor code until each is explicitly
admitted.

Floating point is admitted but is not thereby free: a reachable arithmetic
operation whose result could be an ambiguous NaN payload is a refusal, because
the fragment does not choose among the payloads the IR permits. Bit-preserving
movement and comparison stay eligible. Either way, a corpus of integer fixtures
does not exercise these refusal paths, and a real workload triggers them
immediately.
\end{remark}
\end{step}

% ---------------------------------------------------------------- step 5
\begin{step}{Check that the question is not vacuous}{%
  \clgnew{VAC}{$\Adm,\PairDom_\coal \neq \emptyset$}%
  \clgnote{VAC}{a claim-adequacy gate, not a premise}%
}
\label{step:vacuity}

Here is a way to obtain a proof of \eqref{eq:ni} for any program whatsoever.
Declare preconditions that contradict each other. There is then no admitted
state, hence no admitted low-equal pair, hence the initial state set of the
product is empty, hence the bad state is unreachable. The formula is
unsatisfiable without executing a single instruction, and the checker reports
success.

This is the countermodel that motivates the gate, and it has a subtler second
form: preconditions that admit exactly one state. Now a nonempty admitted set
and a nonempty pair domain both hold, via the pair of that state with itself,
while no secret component actually \emph{varies}. The product is inhabited and
still proves nothing.

So two things must be exhibited, not assumed: at least one admitted execution,
and at least one admitted \emph{coupled low-equal pair} for each coalition. And a
third disposition is needed for the middle case --- a secret component that
provably cannot vary is reported as constrained-or-unexercised rather than
silently contributing to a proof, and the verifier may not infer an
expected-to-vary assertion from a component's name, width, or classification.

\begin{remark}
This gate is not a logical premise of \eqref{eq:ni}; it is a
\emph{claim-adequacy} gate on the report. The distinction matters for the data
structure: both kinds of blocker produce \Unknown{} with a reason, but only
premises appear in the antecedent of the theorem being claimed, so a report that
cannot tell them apart cannot state which theorem it proved.
\end{remark}
\end{step}

% ---------------------------------------------------------------- step 6
\begin{step}{Run the diagnostic layer, and let it fail}{%
  \clgnew{DIAG}{$\DiagOK$}%
  \clgtouch{MAN}%
}
\label{step:diag}

Now the cheap pass from \cref{sec:twosafety}: one \textsf{Low}/\textsf{High}
propagation assigning every site a disposition. Its obligation is only that the
assignment is \emph{total} --- every site classified, nothing silently skipped.

The disposition domain is where the discipline lives:

\begin{center}
\begin{tabular}{@{}ll@{}}
\toprule
not observable & no coalition-visible payload claim at this site \\
statically discharged & closed by a named rule, with premise references \\
definite violation & a violation with a witness recipe \\
\RelReq{} & precision lost; \textbf{the product must decide} \\
\Unknown{} & this \emph{analysis} did not complete: a health failure \\
\bottomrule
\end{tabular}
\end{center}

The last two rows are not two grades of the same thing, and the corpus records
their conflation as a resolved defect. \RelReq{} is a \emph{precision} outcome:
it is this layer working exactly as designed. The diagnostic \Unknown{} is a
\emph{health} failure --- incomplete traversal, a malformed record, a failed
consistency check --- and it is the only diagnostic disposition that blocks
\Proved{}. Keeping them apart is what lets a run say ``I was imprecise here''
without also saying ``I may be broken.''

The last row is narrower than it looks and is worth not misreading. A
\emph{diagnostic} \Unknown{} is reserved for a failure of the diagnostic pass
itself --- it did not achieve its required whole-entry coverage. That is a
different thing from the artifact-level \Unknown{} of \cref{step:normalform} and
\cref{step:verdict}, which reports missing semantics or an absent contract. Two
layers, two \Unknown{}s, and only the second is a verdict about the program.

The fourth value is the one that makes the design work, and the one an
implementer is most likely to try to eliminate. Consider four programs that are
genuinely safe and that this layer cannot see are safe:

\begin{itemize}[nosep,leftmargin=1.4em]
\item a secret written to a slot and then fully overwritten by a public value
      before any read --- needs a strong update, which this layer may not do;
\item a secret at one byte offset with the sink fed from a disjoint offset ---
      needs byte-exact disjointness, and the diagnostic's memory regions are not
      proof certificates;
\item \code{xor} of a value with itself --- needs value congruence, which a
      label lattice collapses away;
\item a branch on a secret whose two edges target the \emph{same} block ---
      structurally lockstep, so disjunct 2 of \cref{sec:twosafety} cannot fire.
\end{itemize}

All four are \RelReq{} here and \Proved{} at the product. None of them is
``silence''; demanding silence is what pushes an implementer toward
\emph{statically discharged} or \emph{not observable}, which \cref{rem:danger}
forbids. There is even a reserved reason code for having taken that shortcut.

The fourth example carries a trap worth its own sentence. The obvious repair ---
``identical successors imply no control leak'' --- is unsound, and
\cref{step:product} exhibits the program it accepts.
\end{step}

% ---------------------------------------------------------------- step 7
\begin{step}{Build the exact product}{%
  \clgnew{PROD}{$\ProdSafe_\coal$}%
  \clgdone{DIAG}%
  \clgtouch{NF}%
}
\label{step:product}

Construct the two-lane product over the whole entry, with the bad-state
disjunction from Part~\ref{part:background}, and ask whether a bad state is
reachable. The diagnostic layer's obligation discharges here: its job was triage,
and the product now carries the burden.

Line~\ref{ln:rel} is the exception that makes the construction correct rather
than merely strict. A release event does not go through the disjuncts at all; it
is routed to the prefix-causal transition of \cref{sec:ledger}, where an
\emph{authorized} release whose two values \emph{differ} is permitted and retires
the pair's obligation. Without that routing every authorized declassification
would register as a leak, and the checker would reject every useful program.

\begin{algorithm}[H]
\caption{Bad-state predicate for coalition $\coal$ at aligned step $k$. Applies
\emph{outside} the release transition of \cref{sec:ledger}; a release event is
routed there instead, and an authorized release with differing values is
permitted rather than bad.}
\label{alg:bad}
\begin{algorithmic}[1]
\State \label{ln:rel} \textbf{if} this step is a release event \textbf{then return} \textsc{ledgerStep}
\State \label{ln:step} \textbf{if} exactly one lane can take a step \textbf{then return} \textsc{bad}
\State \label{ln:loc}  \textbf{if} the lanes' next control locations differ \textbf{then return} \textsc{bad}
\State \label{ln:status} \textbf{if} their termination or failure classes differ \textbf{then return} \textsc{bad}
\State \label{ln:align} \textbf{if} the emitted event's site or occurrence index differs \textbf{then return} \textsc{bad}
\State \label{ln:word} \textbf{if} $\Proj_\coal(\text{event}_0) \neq \Proj_\coal(\text{event}_1)$ \textbf{then return} \textsc{bad}
\State \textbf{return} \textsc{ok}
\end{algorithmic}
\end{algorithm}

Three places where an independent implementation plausibly diverges, and where
divergence is unsound rather than merely different:

\textbf{Line~\ref{ln:loc} keys on successor identity, not on the condition's
label.} A branch on a secret into one block has a singleton successor set and
cannot fire this line, which is exactly why the fourth example in
\cref{step:diag} is safe. But now consider a branch on a secret into one block
whose two \emph{edges carry different operands} --- in the LLVM dialect a
$\phi$ is a block argument, so the merged value arrives as an argument of the
target block with no operand edge from the branch. Lines~\ref{ln:step}
through~\ref{ln:align} all pass. The leak lands only on Line~\ref{ln:word}.

That is the program the repair rejected in \cref{step:diag} would accept, and it
is not hypothetical: canonicalizing an ordinary two-armed diamond in
\code{mlir-opt} produces precisely this shape. It also refutes a more general
principle --- that an ordinary SSA slice is closed around a $\phi$ --- since a
dependence relation built over operand edges never sees the gating fact.

\textbf{Line~\ref{ln:align} is global, not per-site.} Comparing occurrence counts
\emph{per event template} is not the same as comparing the global interleaving.
Two lanes can emit the same events, the same number of times, with identical
payloads, in the opposite order. Every per-template row agrees; the traces
differ. So the projected trace must be an ordered sequence with pairwise index
alignment and a first-divergence position --- never a multiset, and never a
map from site to last value.

\textbf{Lines~\ref{ln:step}--\ref{ln:align} do not weaken under projection.}
They are structural, so a coalition that cannot read a payload still
distinguishes these executions. An implementation that applies the coalition
projection before evaluating them will wrongly equate lanes that differ in
control location, count, or termination.

\begin{remark}[What may not be filtered]
An execution that exhausts an unrolling bound, faults, or risks undefined
behaviour is \emph{retained}, not deleted. Deleting it silently narrows the proof
domain --- the same vacuity failure as \cref{step:vacuity} arriving by a
different route. Bound adequacy is a universal obligation discharged separately,
not a filter applied inside the admitted set.
\end{remark}
\end{step}

% ---------------------------------------------------------------- step 8
\begin{step}{Prove, refute, or refuse}{%
  \clgdone{PROD}%
  \clgdone{VAC}%
  \clgnew{WIT}{$\ReplayCov_\coal(w)$}%
  \clgdone{WIT}%
}
\label{step:verdict}

The product query returns one of three things, and the third is not a failure of
the tool.

\textbf{Unreachable} $\Rightarrow$ \Proved{}, provided the adequacy gate of
\cref{step:vacuity} holds. That proviso is why both obligations discharge in the
same step.

\textbf{Reachable, with a witness} $\Rightarrow$ \Ctrex{}, but only if the
witness \emph{replays}. A model returned by a solver is a candidate, not a
finding: it must be re-executed under the exact semantics and shown to reach the
bad state, with its prefix covered. Failed replay or missing coverage is
\Unknown{} and tool inconsistency --- never a counterexample. This is the
obligation that opens and closes inside this step.

\textbf{Neither} $\Rightarrow$ \Unknown{}, with the blocker set named.

Two rules about aggregation, both of which are easy to get wrong and neither of
which is cosmetic.

\emph{A replayed witness outranks an open universal obligation.} If a witness
reaches a bad state and replays, the verdict is \Ctrex{} regardless of an
unfinished bound-adequacy proof, an unfinished definedness proof, or another
solver timeout elsewhere --- and regardless of whether a solver, a fuzzer, a
differential search, or a person found it.

Source-independence is a requirement on the \emph{verdict}, not on the record.
The result must be identical from every discovery route, and no step of the
argument may consult provenance. The record itself is the opposite of anonymous:
it must \textbf{retain} the discovery source and the seed or model, precisely so
the candidate can be re-executed and the replay audited. Confusing the two ---
stripping provenance in the name of source-independence --- destroys
reproducibility while buying nothing, since it is the argument that must ignore
the source, not the file.

\emph{Otherwise every blocker is preserved.} With one blocker the result is
\Unknown{} carrying it. With several it is \Unknown{} carrying a compound reason
over the \emph{complete ordered} blocker set. A flat enumeration of reason codes
cannot represent that, which makes the compound constructor a genuine
requirement on the record type rather than a convenience.

\begin{remark}
Note the shape of the record this forces. \Unsafe{} carries a replayable
witness; \Unknown{} carries an ordered blocker set; \Proved{} carries neither.
A single ``reason'' string field cannot hold all three, and a design that tries
will end up attaching reason codes to proofs, which are reason-free.
\end{remark}
\end{step}

% ---------------------------------------------------------------- step 9
\begin{step}{Close the deployment gap, or admit it is open}{%
  \clgnew{DEP}{$\PairedRef_\coal$}%
  \clgnote{DEP}{leaving this open is what \Conditional{} means}%
}
\label{step:deployment}

Everything so far concerns a frozen artifact under a declared machine model. Two
gaps remain between that and a running system, and they are of different kinds.

\Lv{3} is \emph{attribution}: relating source, imported IR, and the final target
artifact, so a violation or a repair can be ascribed to a specific compiler
boundary. This is where the classic ``the compiler broke my constant-time code''
findings live --- a conditional move lowered to a branch on a target without
conditional moves, or a division strength-reduced into a multiply-high so the
operation you were reasoning about is not the operation that executes.

\Lv{4} is \emph{facts outside the modeled semantics}: real helper latency,
compressor behaviour, hardware isolation between tenants, sanitizer sufficiency,
certificate soundness, and whether the backend preserved the observable trace at
all. These are not provable from the IR. They are empirical or hardware facts.

The honest handling is the fourth outcome. A repair that is correct in the model
but depends on an undischarged external fact is \Conditional{}, and must name
every assumption it rests on. An \Unsafe{} or \Verified{} result may name none.

\begin{remark}[Why this obligation is per-coalition and separate]
Deployment status is an independent axis from model status. A model-level
counterexample cannot be repaired by deployment evidence, and deployment evidence
cannot upgrade an \Unknown{} model result. Collapsing the two into one verdict
token is the most common structural error, and it shows up concretely: an
obligation like ``the backend preserved the trace'' and one like ``this helper's
latency is operand-independent'' belong to different axes and should not share a
field.
\end{remark}
\end{step}

% ---------------------------------------------------------------- step 10
\begin{step}{Close the arc}{%
  \clgdone{DEP}%
  \clgdone{NF}%
  \clgdone{MAN}%
  \clgdone{BIND}%
  \clgdone{NI}%
}
\label{step:close}

Discharge the artifact-level gates and the goal. The manifest was total, the
artifact conformed, the results are bound to those exact bytes and that exact
pipeline, and the deployment obligations of \cref{step:deployment} were
\emph{discharged}. The product was unreachable under a non-vacuous admitted
domain, so \eqref{eq:ni} holds for every derived coalition at the claimed
boundary.

This is the \Verified{} branch, and it is the only branch in which the arc closes.
Had the deployment obligations merely been \emph{named} rather than discharged,
the outcome would be \Conditional{} --- and the ledger would end with that row
still open, which is exactly the invariant from Part~\ref{part:orientation}:
\Verified{} may carry no obligations, so a non-empty ledger is not a proof but a
different verdict. An empty ledger is the picture of one specific outcome, not of
the analysis having run.

What has \emph{not} been established is worth saying explicitly, because it is
the difference between an honest result and an overclaim:

\begin{itemize}[nosep,leftmargin=1.4em]
\item nothing about constructs outside the fragment of \cref{step:normalform};
\item nothing about a different artifact, pipeline, or target profile;
\item nothing about real latency, hardware isolation, or speculation beyond what
      \cref{step:deployment} explicitly named and discharged;
\item nothing about coalitions of principals absent from the manifest, since the
      closure was taken over the declared set.
\end{itemize}

Every one of those is a boundary, not a bug. The value of the four-outcome
contract is that it can state them.
\end{step}

\clgboard

\noindent
The ledger is empty. Most middle steps traded an obligation for a successor that
was easier to attack; a few only opened one, because some preconditions have to
be stated before anything can be traded against them. The two that stayed open
longest were the least glamorous: that the policy was total, and that the answer
was about the artifact we actually shipped.
